\Needspace{10\baselineskip}
\section{Comparing local regularizations}

To glue local convolutions, we need them to agree up to an error tending uniformly to zero on overlaps. A bounded plurisubharmonic function need not be continuous, so this comparison cannot come from uniform convergence to the original function. Instead, logarithmic convexity of its suprema on balls controls the change of scale and the change of coordinates.

\Needspace{9\baselineskip}
Let $U\subset\mathbb C^n$ be a domain and $u$ a bounded plurisubharmonic function on $U$. On a fixed relatively compact subset, for small $\delta>0$ put
\[
 S_\delta u(z)=\sup_{|\zeta-z|\leq\delta}u(\zeta),\qquad
 A_\delta u(z)=\fint_{|\zeta-z|=\delta}u(\zeta)\,\mathrm{d}\sigma(\zeta),
 \qquad u_\delta=u*\rho_\delta,
\]
where $\mathrm{d}\sigma$ is Euclidean surface measure and $\rho$ is a smooth nonnegative radial mollifier of integral one, supported in the unit ball, with $\rho_\delta(\zeta)=\delta^{-2n}\rho(\zeta/\delta)$. A fixed larger neighborhood contains all balls below, and $B$ denotes a bound for the oscillation of $u$ there.

\Needspace{16\baselineskip}
\begin{lemma}\label{lem:bounded-local-regularization}
For each fixed $A>1$,
\[
 S_{A\delta}u-S_\delta u\longrightarrow0,
 \qquad S_\delta u-u_\delta\longrightarrow0
\]
uniformly on compact subsets as $\delta\downarrow0$.
\end{lemma}
\begin{proof}
The maximum of a plurisubharmonic function on concentric balls is convex as a function of the logarithm of the radius. This follows by restricting to each complex line through the center and comparing on an annulus with the harmonic function affine in $\log|\zeta|$. For a fixed admissible radius $r>A\delta$, convexity gives
\begin{equation}\label{eq:bounded-log-radius}
 0\leq S_{A\delta}u-S_\delta u
 \leq\frac{B\log A}{\log(r/\delta)}.
\end{equation}
This proves the first assertion.

We compare the maximum with the spherical mean using the Poisson kernel in real dimension $2n$. For a point at distance at most $\delta/2$ from the center, the normalized Poisson kernel of the ball of radius $\delta$ is bounded below by a positive dimensional constant $c_n$. Since $u\leq S_\delta u$, the submean comparison with its harmonic extension gives
\[
 S_\delta u-S_{\delta/2}u
 \geq c_n(S_\delta u-A_\delta u).
\]
These inequalities are valid for bounded plurisubharmonic functions by local smoothing, followed by a decreasing limit. Equation \eqref{eq:bounded-log-radius} shows that $S_\delta u-A_\delta u$ tends uniformly to zero.

For $0<t\leq1$, the spherical mean is increasing with radius, and
\[
 0\leq A_\delta u-A_{t\delta}u
 \leq (S_\delta u-S_{t\delta}u)
                +(S_{t\delta}u-A_{t\delta}u).
\]
The first term is at most
$B\log(1/t)/\log(r/(t\delta))$ by logarithmic convexity of $S$, and the second is at most $C_nB/\log(r/(t\delta))$ by the preceding estimate. Write the radial convolution as an average of $A_{t\delta}u$ with the probability density
$\widehat\rho(t)=|\mathbb S^{2n-1}|t^{2n-1}\rho(t)$. Since
$\int_0^1(1+\log(1/t))\widehat\rho(t)\,\mathrm{d}t<\infty$, integration gives
\[
 0\leq A_\delta u-u_\delta
 \leq\frac{C_{n,\rho}B}{\log(r/\delta)}.
\]
Combining the two comparisons proves the second assertion.
\end{proof}

\Needspace{9\baselineskip}
\begin{corollary}\label{cor:coordinate-convolution}
Let $U,V\subset\mathbb C^n$ be domains and let $F:U\to V$ be biholomorphic and $u$ bounded and plurisubharmonic on $U$. The difference
\[
 u_\delta-\bigl((u\circ F^{-1})*\rho_\delta\bigr)\circ F
\]
tends locally uniformly to zero as $\delta\downarrow0$.
\end{corollary}
\begin{proof}
On a compact subset, $F$ and its inverse have uniformly bounded first derivatives. Thus, for a fixed $A>1$, the inverse image of a radius-$\delta$ ball about $F(z)$ lies between balls about $z$ of radii $\delta/A$ and $A\delta$. The two supremum regularizations are therefore bounded by one another at these enlarged radii. The first assertion of \cref{lem:bounded-local-regularization} makes their difference uniformly small. Its second assertion compares each of the two convolutions with the corresponding supremum regularization.
\end{proof}

\Needspace{12\baselineskip}
\section{Gluing with a small curvature error}

The overlap estimate permits gluing by a smooth maximum. Small boundary corrections make each local branch irrelevant near the edge of its chart, while convexity of the maximum preserves a common lower curvature bound. Use $\omega$ to measure the curvature error.

For an integer $N\geq1$, choose independent real random variables $Z_1,\ldots,Z_N$ with a common smooth even probability density supported in $[-1,1]$, and define
\[
 M_\tau(t_1,\ldots,t_N)
     =\mathbb E\max_i(t_i+\tau Z_i)\quad(\tau>0).
\]
Convolution makes $M_\tau$ smooth. It is convex, increasing in every argument, and, with $\mathbf1=(1,\ldots,1)$, satisfies
\[
 \max_i t_i\leq M_\tau(t)\leq\max_i t_i+\tau,
 \qquad M_\tau(t+c\mathbf1)=M_\tau(t)+c
\]
for $t\in\mathbb R^N$ and $c\in\mathbb R$. If $t_i<\max_{j\ne i}t_j-2\tau$, the $i$th argument can be removed without changing the value. Also $M_\tau(t)$ is nondecreasing in $\tau$: it is a convex even function of $\tau$ because the distribution of $Z$ is symmetric.

The second identity expresses translation invariance. If smooth real functions $f_i$, $1\leq i\leq N$, satisfy $\ddc f_i\geq-\gamma$ for a common smooth real $(1,1)$-form $\gamma$, then
\begin{equation}\label{eq:smooth-max-curvature}
 \ddc M_\tau(f_1,\ldots,f_N)\geq-\gamma.
\end{equation}
Indeed, the first-derivative coefficients are nonnegative and sum to one, by the translation identity, while the real Hessian of the convex function $M_\tau$ contributes a semipositive $(1,1)$-form in the chain rule.

\Needspace{8\baselineskip}
\begin{proposition}\label{prop:bounded-almost-regularization}
Let $\theta$ be smooth and real, $\omega$ Hermitian, and
$u\in\PSH(X,\theta)\cap L^\infty(X)$. There are smooth functions $v_j$ and positive numbers $\varepsilon_j\downarrow0$ such that
\[
 v_j\downarrow u,\qquad
 \theta+\varepsilon_j\omega+\ddc v_j\geq0.
\]
\end{proposition}
\Needspace{9\baselineskip}
\begin{proof}
Fix $\epsilon>0$. By continuity of $\theta$, choose finitely many nested coordinate neighborhoods
$V_i\Subset W_i\Subset U_i$, with the $V_i$ covering $X$, and smooth quadratic functions $r_i$ such that
\[
 0\leq\ddc r_i-\theta\leq\epsilon\omega\quad\hbox{on }U_i.
\]
To make this choice, freeze the coefficients of $\theta$ at the center, add a small positive constant form, and shrink the coordinate neighborhood. Put $p_i=u+r_i$, a bounded plurisubharmonic function, and let $p_{i,\delta}$ be its convolution in $U_i$.

On an overlap, express both $p_i$ and $p_j$ in the $j$th coordinates. Their difference is the smooth function $r_i-r_j$, so the difference of their convolutions tends uniformly to $r_i-r_j$. By \cref{cor:coordinate-convolution}, changing the coordinates used to convolve $p_i$ introduces another uniformly vanishing error. Thus, on overlaps of the closures of the $W_i$,
\begin{equation}\label{eq:overlap-regularization}
 (p_{i,\delta}-r_i)-(p_{j,\delta}-r_j)\longrightarrow0
 \quad\hbox{uniformly}.
\end{equation}
Choose $\eta_i\in C^\infty(U_i,[-1,0])$, zero on $V_i$ and equal to $-1$ near the boundary of $W_i$, and $C\geq1$ with $\ddc\eta_i\geq-C\omega$. The cover, the functions $r_i,\eta_i$ and the constant $C$ may depend on $\epsilon$; they are fixed while $\delta$ tends to zero. Set
\[
 f_{i,\delta}=p_{i,\delta}-r_i+(\epsilon/C)\eta_i
 \quad\hbox{on }W_i.
\]
Each satisfies $\ddc f_{i,\delta}\geq-\theta-2\epsilon\omega$. Near a boundary point of $W_i$, some $V_j$ contains that point. Equation \eqref{eq:overlap-regularization} implies, for all small $\delta$,
\[
 f_{i,\delta}<f_{j,\delta}-\epsilon/(2C)
\]
on a neighborhood of that boundary point. Finiteness of the cover makes this separation uniform.

Take $\tau(\delta)\downarrow0$ as $\delta\downarrow0$, with $\tau(\delta)<\epsilon/(8C)$, and define on each point the smooth maximum of the branches whose $W_i$ contain it. A branch disappears only near the boundary of its $W_i$, where it is below another branch by more than $2\tau$. The removal property of $M_\tau$ therefore proves that these formulas define a global smooth function $\psi_{\epsilon,\delta}$. Equation \eqref{eq:smooth-max-curvature} gives
\[
 \theta+2\epsilon\omega+\ddc\psi_{\epsilon,\delta}\geq0.
\]
At every point there is a branch with $\eta_i=0$, and $p_{i,\delta}\geq p_i$. Hence $\psi_{\epsilon,\delta}\geq u$. For fixed $\epsilon$, the local convolutions decrease with $\delta$, and so does the smooth maximum when both its arguments and $\tau$ decrease. Its pointwise limit is
\[
 \max_{i:x\in W_i}\bigl(u(x)+(\epsilon/C)\eta_i(x)\bigr)=u(x).
\]
Thus $\psi_{\epsilon,\delta}\downarrow u$ as $\delta\downarrow0$, also in $L^1$.

Choose a decreasing diagonal as follows. Take $\epsilon_j\downarrow0$ and put
$v_j=\psi_{\epsilon_j,\delta_j}+2^{-j}$. Choose $\delta_j$ so small that
\[
 \|\psi_{\epsilon_j,\delta_j}-u\|_{L^1}\leq2^{-j},
 \qquad v_j<v_{j-1}\quad(j\geq2).
\]
The second choice is possible because $v_{j-1}\geq u+2^{-j+1}$. More explicitly, the continuous functions
$\psi_{\epsilon_j,\delta}+2^{-j}-v_{j-1}$ decrease pointwise to a function bounded above by $-2^{-j}$. Their nonnegative level sets are decreasing compact sets with empty intersection, so one of them is empty. The resulting sequence decreases, converges to $u$ in $L^1$, and stays above $u$. Its pointwise limit equals $u$ by uniqueness of canonical quasi-plurisubharmonic representatives. Taking $\varepsilon_j=2\epsilon_j$ proves the proposition.
\end{proof}

\Needspace{16\baselineskip}
\section{Regularization in a fixed Hermitian background}

For $\theta=\omega$, rescaling the approximants absorbs the curvature error into the positive background. Subtracting a constant first makes the approximants negative, so that this rescaling preserves their decreasing order.

\begin{theorem}\label{thm:bounded-exact-regularization}\index{Regularization!bounded potentials}
For $u\in\PSH(X,\omega)\cap L^\infty(X)$ there are smooth functions $u_j\downarrow u$ such that
\[
 \omega+\ddc u_j>0.
\]
The functions $u_j$ are uniformly bounded.
\end{theorem}
\begin{proof}
Subtract a constant so that $u\leq-2$. Apply \cref{prop:bounded-almost-regularization} with $\theta=\omega$, obtaining smooth $v_j\downarrow u$ and positive $\varepsilon_j\downarrow0$ with $(1+\varepsilon_j)\omega+\ddc v_j\geq0$. The compact sets $\{v_j\geq-1\}$ decrease to the empty set, since $u\leq-2$. Discarding finitely many terms therefore gives $v_j<0$ everywhere. Write
\[
 u_j=\frac{v_j}{1+2\varepsilon_j}.
\]
Then
\[
 \omega+\ddc u_j
 \geq\frac{\varepsilon_j}{1+2\varepsilon_j}\omega>0.
\]
Because $v_j$ is negative and decreases while the positive denominator decreases to one, $u_j$ also decreases and converges to $u$. It lies between $u$ and its first smooth member, which gives uniform boundedness. Restore the subtracted constant.
\end{proof}

This regularization supplies the approximation used for products and mass in \cref{chap:products}. For an arbitrary smooth $\theta$, \cref{prop:bounded-almost-regularization} retains a vanishing curvature error.
